\section*{第 \ref{ch:b3:advanced-nmr} 章——核磁共振进阶：脉冲、弛豫与二维谱}

\begin{solution}{exo:b3:advanced-nmr:1}
$\nu_0 = (\gamma/2\pi)B_0$：\ce{^1H} \qty{600.4}{MHz}，\ce{^{13}C} \qty{151.0}{MHz}，\ce{^{19}F}
\qty{565.1}{MHz}。
\end{solution}

\begin{solution}{exo:b3:advanced-nmr:2}
$h\nu_0/2kT = 6.626\times10^{-34} \times 400.2\times10^6/(2 \times 1.381\times10^{-23} \times 300)
= \num{3.2e-5}$。在 \qty{4}{K} 时该值为原来的 75 倍，即 \num{2.4e-3}（近似
仍然成立）。
\end{solution}

\begin{solution}{exo:b3:advanced-nmr:3}
$\theta = \gamma B_1t_p = \pi/2$ 对应 \qty{10}{\micro s}：$\gamma B_1/2\pi = 1/(4 \times
\qty{10}{\micro s}) = \qty{25}{kHz}$。$180^\circ$ 脉冲持续 \qty{20}{\micro s}；\qty{5}{\micro s}
给出 $45^\circ$。
\end{solution}

\begin{solution}{exo:b3:advanced-nmr:4}
$T_2 = 1/(\pi \times 3.2) = \qty{0.10}{s}$.
\end{solution}

\begin{solution}{exo:b3:advanced-nmr:5}
$(10.71/42.58)^3 = 0.0159$，乘以 0.011 后得 \num{1.75e-4}，约为质子信号的
1/5700。由于 S/N $\propto\sqrt N$，要使 S/N 相同，扫描次数需增加
$5700^2 \approx \num{3e7}$ 倍，实际不可行，这正是 ${}^{13}$C 谱采用更浓样品、
去偶和极化转移的原因。
\end{solution}

\begin{solution}{exo:b3:advanced-nmr:6}
$T_1 = 1.39/\ln2 = \qty{2.0}{s}$；循环延迟 $\ge 5T_1 = \qty{10}{s}$。
\end{solution}

\begin{solution}{exo:b3:advanced-nmr:7}
COSY 谱中只有一个交叉峰，位于 \ce{CH2} 四重峰与
乙基 \ce{CH3} 三重峰之间；\ce{CH3CO} 单峰与任何质子都不耦合。由单峰甲基出发的 HMBC 相关为：
羰基碳（两键）和 \ce{CH2} 碳（三键）。
\end{solution}

\begin{solution}{exo:b3:advanced-nmr:8}
$k_c = \pi \times 50/\sqrt2 = \qty{111}{s^{-1}}$. $\Delta G^\ddagger = 8.314 \times 330 \times
\ln(6.88\times10^{12}/111) = \qty{68}{kJ/mol}$.
\end{solution}

\begin{solution}{exo:b3:advanced-nmr:9}
核欧沃豪斯效应 $\propto r^{-6}$：$r_b = 0.30 \times 4^{-1/6} = \qty{0.24}{nm}$。
\end{solution}

\begin{solution}{exo:b3:advanced-nmr:10}
在共振的旋转坐标系中，只剩 $\mathbf B_1 = B_1\mathbf e_{x'}$：
$\dot M_{x'} = 0$，$\dot M_{y'} = \gamma B_1M_z$，$\dot M_z = -\gamma B_1M_{y'}$（
$\gamma\mathbf M\times\mathbf B_1$ 的分量）。于是 $M_{x'}$ 为常数，$(M_{y'}, M_z)$ 以角速率
$\gamma B_1$ 转动：经 $t_p$ 后，转过角 $\gamma B_1t_p$，绕 $x'$ 轴。
\end{solution}

\begin{solution}{exo:b3:advanced-nmr:11}
处于稍强磁场中的自旋，在
$\tau$ 期间以恒定的附加速率获得相位；$180^\circ$ 脉冲使其相位反转，并在
第二个 $\tau$ 期间损失完全相同的相位：所有此类自旋在 $2\tau$ 时重新同相。随机的自旋–自旋
相互作用随时间涨落，在第二段时间内不再重复，因此其
失相不能被消除。回波幅度按真正的 $T_2$ 衰减。
\end{solution}

\begin{solution}{exo:b3:advanced-nmr:12}
乙酸乙酯，\ce{CH3COOCH2CH3}：\ce{OCH2} 四重峰（4.12），\ce{CH3CO} 单峰（2.05），
\ce{CH3} 三重峰；四重峰的质子与羰基碳相隔三个键，
经由酯氧相连。若为丙酸甲酯，\ce{CH3CH2COOCH3}，则其
四重峰位于 2.3 ppm 附近（在 HSQC 中对应碳约 27 ppm，而非 60），单峰
\ce{OCH3} 位于 3.7 ppm 附近，其到羰基的 HMBC 峰穿过氧原子。
\end{solution}

\begin{solution}{pb:b3:advanced-nmr:1}
\textbf{1.} $(2 \times 6 + 2 - 12)/2 = 1$：一个 C=O。
\textbf{2.} 173.6 ppm：酯羰基；60.1 ppm：\ce{OCH2}。
\textbf{3.} 五个；DEPT-135：\ce{CH2} 为负峰（60.1、36.3、18.6），\ce{CH3} 为正峰（14.3、
13.7），羰基碳不出现。
\textbf{4.} \qty{400.2}{MHz} 与 \qty{100.7}{MHz}。
\textbf{5.} $3J = \qty{21.6}{Hz} = \qty{0.054}{ppm}$.
\textbf{6.} 多重性表明每个片段内部哪些基团相邻，
但无法说明各片段如何通过氧和羰基相连。
\textbf{7.} 4.13--1.25；2.28--1.66；1.66--0.95（以及对称的那些）。
\textbf{8.} \ce{OCH2CH3}（4.13、1.25）与 \ce{CH2CH2CH3}（2.28、1.66、0.95）。
\textbf{9.} \ce{OCH2} 与 \ce{CH2C=O} 质子之间隔着氧原子和
羰基碳，相隔五个键，无法分辨出耦合。
\textbf{10.} 1.66 ppm，位于 \ce{CH2} 与 \ce{CH3} 之间，与 2 + 3 = 5 个
$J$ 相近的质子耦合，呈六重峰（即一种多重峰）。
\textbf{11.} 两个甲基相互耦合，因而两个 \ce{CH3} 位于相邻碳上（\ce{CH3-CH3} 单元）：此处不可能。
\textbf{12.} \ce{-O-CH2-CH3}、\ce{-CH2-CH2-CH3}，以及 \ce{C=O}。
\textbf{13.} 4.13/60.1；2.28/36.3；1.66/18.6；1.25/14.3；0.95/13.7。
\textbf{14.} 14.3（乙基 \ce{CH3}）与 13.7 ppm（丁酰基 \ce{CH3}），仅相差 \qty{0.6}{ppm}。
\textbf{15.} 4.13 ppm 到 173.6 ppm：H--C--O--C，三个键。
\textbf{16.} 2.28 ppm（两键）和 1.66 ppm（三键）与 173.6 ppm 相关。
\textbf{17.} \ce{CH3CH2CH2C(=O)OCH2CH3}，即丁酸乙酯。丙酸丙酯会在 4.0 ppm 附近给出
\ce{OCH2} 三重峰（而非四重峰），并在 2.3 ppm 附近出现四重峰。
\textbf{18.} 四键耦合通常小于 1 Hz；HMBC 延迟设定为
5--10 Hz，因而不会选中它们。
\textbf{19.} 不同碳的弛豫速率不同，去偶引起的核
欧沃豪斯增强也不同。
\textbf{20.} $1 - \eu^{-5/20} = 0.22$：羰基碳信号偏弱，约低 4 倍。
\textbf{21.} $1 - \eu^{-5} = 0.993$.
\textbf{22.} 质子去偶产生的核欧沃豪斯效应因碳而异，
可用反门控去偶予以消除（去偶器仅在采集期间开启）。
\textbf{23.} $256 \times \qty{100}{s} = \qty{25600}{s}$，约 7 小时。
\textbf{24.} \textbf{羰基碳的 $5T_1$：两次扫描间隔 \qty{100}{s}。}
\end{solution}
